% Source: Weinberg GR, physical PDF pp. 357--360
%         (printed pp. 335--338).
% Coverage: Section 11.7, Time-Dependent Spherically Symmetric Fields;
%           equations (11.7.1)--(11.7.9).
% Figures/tables/footnotes: no figures or tables; reference callout 30.
% Status: source-reviewed and compile-clean.
% Uncertainties: the curvature-sign conversion is documented at
%                Eqs. (11.7.4)--(11.7.7); none unresolved.

\subsection{Time-Dependent Spherically Symmetric Fields}\label{sec:11.7}

We now turn to the problems of stellar dynamics, and begin by writing down the
metric and Ricci tensor for a spherically symmetric but time-dependent system.
Spherical symmetry requires the invariant interval \(\dd s^2\) to depend only on
the rotational invariants
\[
  t,\quad \dd t,\quad r,\quad
  \mathbf{x}\cdot \dd\mathbf{x}=r\,\dd r,\quad
  \dd\mathbf{x}^2=\dd r^2+r^2(\dd\theta^2+\sin^2\theta\,\dd\varphi^2),
\]
so, in the binding signature, it can be written
\[
  \begin{aligned}
  \dd s^2={}&
  -C(r,t)\,\dd t^2+D(r,t)\,\dd r^2+2E(r,t)\,\dd r\,\dd t\\
  &+F(r,t)r^2(\dd\theta^2+\sin^2\theta\,\dd\varphi^2).
  \end{aligned}
\]
The function \(F\) can be removed by defining a new radial variable
\[
  r'\equiv rF^{1/2}(r,t).
\]
The metric will then be of the same form, but with new functions
\(C',D',E'\) in place of \(C,D,E\), and of course with \(r'\) in place of
\(r\) and no factor \(F\). Dropping primes, we have then
\[
  \begin{aligned}
  \dd s^2={}&
  -C(r,t)\,\dd t^2+D(r,t)\,\dd r^2+2E(r,t)\,\dd r\,\dd t\\
  &+r^2(\dd\theta^2+\sin^2\theta\,\dd\varphi^2).
  \end{aligned}
\]
We next remove \(E\), by defining a new time
\[
  \dd t'=\eta(r,t)[C(r,t)\,\dd t-E(r,t)\,\dd r],
\]
where \(\eta\) is an integrating factor defined to make the right-hand side a
perfect differential, that is, so that
\[
  \partial_r[\eta(r,t)C(r,t)]
  =
  -\partial_t[\eta(r,t)E(r,t)].
\]
(This equation can be solved by treating it as an initial-value problem:
given \(\eta(r,t_0)\) for all \(r\), we can solve for
\(\partial_t\eta(r,t)|_{t=t_0}\) and thus determine
\(\eta(r,t_0+\dd t)\) for all \(r\).) The interval is then
\[
  \dd s^2
  =
  -\eta^{-2}C^{-1}\,\dd t'^2
  +(D+C^{-1}E^2)\,\dd r^2
  +r^2(\dd\theta^2+\sin^2\theta\,\dd\varphi^2),
\]
or, introducing new functions \(A\) and \(B\) in place of
\(D+C^{-1}E^2\) and \(\eta^{-2}C^{-1}\), and dropping the prime on \(t\),
\begin{equation}
  \dd s^2
  =
  -B(r,t)\,\dd t^2+A(r,t)\,\dd r^2
  +r^2(\dd\theta^2+\sin^2\theta\,\dd\varphi^2).
  \tag{11.7.1}\label{eq:11.7.1}
\end{equation}
Thus we can use the metric in its familiar ``standard'' form, the only new
feature being that \(A\) and \(B\) now depend on \(t\) as well as \(r\).

The nonvanishing elements of the metric tensor and its inverse are
\begin{equation}
  \begin{gathered}
  g_{tt}=-B,\qquad
  g_{rr}=A,\qquad
  g_{\theta\theta}=r^2,\qquad
  g_{\varphi\varphi}=r^2\sin^2\theta,\\
  g^{tt}=-B^{-1},\qquad
  g^{rr}=A^{-1},\qquad
  g^{\theta\theta}=r^{-2},\qquad
  g^{\varphi\varphi}=r^{-2}(\sin\theta)^{-2}.
  \end{gathered}
  \tag{11.7.2}\label{eq:11.7.2}
\end{equation}
It follows that the nonvanishing elements of the affine connection are
\begin{equation}
  \begin{gathered}
  \Gamma^r{}_{rr}=\frac{A'}{2A},\qquad
  \Gamma^r{}_{\theta\theta}=-\frac rA,\qquad
  \Gamma^r{}_{\varphi\varphi}=-\frac{r\sin^2\theta}{A},\qquad
  \Gamma^r{}_{tt}=\frac{B'}{2A},\\
  \Gamma^r{}_{rt}=\Gamma^r{}_{tr}=\frac{\dot A}{2A},\qquad
  \Gamma^\theta{}_{r\theta}=\Gamma^\theta{}_{\theta r}=\frac1r,\qquad
  \Gamma^\theta{}_{\varphi\varphi}=-\sin\theta\cos\theta,\\
  \Gamma^\varphi{}_{r\varphi}=\Gamma^\varphi{}_{\varphi r}=\frac1r,\qquad
  \Gamma^\varphi{}_{\theta\varphi}
  =\Gamma^\varphi{}_{\varphi\theta}=\cot\theta,\\
  \Gamma^t{}_{rr}=\frac{\dot A}{2B},\qquad
  \Gamma^t{}_{tt}=\frac{\dot B}{2B},\qquad
  \Gamma^t{}_{rt}=\Gamma^t{}_{tr}=\frac{B'}{2B}.
  \end{gathered}
  \tag{11.7.3}\label{eq:11.7.3}
\end{equation}
(A prime or a dot now denotes \(\partial_r\) or \(\partial_t\),
respectively.) From the definition of the Ricci tensor in the binding
curvature convention we obtain the independent nonzero components:
% MODERNIZATION CHECK: every component in (11.7.4)--(11.7.7) is the
% negative of Weinberg's printed Ricci component.
\begin{align}
  R_{rr}
  ={}&
  -\frac{B''}{2B}
  +\frac{B'^2}{4B^2}
  +\frac{A'B'}{4AB}
  +\frac{A'}{rA}
  +\frac{\ddot A}{2B}
  -\frac{\dot A\dot B}{4B^2}
  -\frac{\dot A^2}{4AB},
  \tag{11.7.4}\label{eq:11.7.4}\\
  R_{\theta\theta}
  ={}&
  1-\frac1A+\frac{rA'}{2A^2}-\frac{rB'}{2AB},
  \tag{11.7.5}\label{eq:11.7.5}\\
  R_{tt}
  ={}&
  \frac{B''}{2A}
  -\frac{B'A'}{4A^2}
  +\frac{B'}{Ar}
  -\frac{B'^2}{4AB}
  -\frac{\ddot A}{2A}
  +\frac{\dot A^2}{4A^2}
  +\frac{\dot A\dot B}{4AB},
  \tag{11.7.6}\label{eq:11.7.6}\\
  R_{tr}
  ={}&
  \frac{\dot A}{Ar}.
  \tag{11.7.7}\label{eq:11.7.7}
\end{align}
Also, it follows from the spherical symmetry of the metric that
\begin{equation}
  R_{\varphi\varphi}=\sin^2\theta\,R_{\theta\theta},
  \tag{11.7.8}\label{eq:11.7.8}
\end{equation}
and
\begin{equation}
  R_{r\theta}=R_{r\varphi}=R_{\theta\varphi}
  =R_{\theta t}=R_{\varphi t}=0.
  \tag{11.7.9}\label{eq:11.7.9}
\end{equation}

As a simple but important application of these results, let us consider a
spherically symmetric but not necessarily static field in \emph{empty space},
where the field equations read \(R_{\mu\nu}=0\). According to
Eq.~\eqref{eq:11.7.7}, the field equation \(R_{tr}=0\) just tells us that
\(A\) is time-independent:
\[
  \dot A=0.
\]
Inspection of Eqs.~\eqref{eq:11.7.4}--\eqref{eq:11.7.6} then shows that
\emph{all} time derivatives drop out of the field equations, and they become
identical with the equations for a static isotropic gravitational field in
empty space. [See Eq.~\eqref{eq:8.1.13}.] We can then repeat the arguments
of Section~\ref{sec:8.2}; the vanishing of \(R_{rr}\) and \(R_{tt}\) gives
\[
  (AB)'=0,
\]
and the vanishing of \(R_{\theta\theta}\) gives
\[
  \left(\frac rA\right)'=1.
\]
Since \(A\) is time-independent, the general solution is
\[
  A=\left(1-\frac{2MG}{r}\right)^{-1},
  \qquad
  B=f(t)\left(1-\frac{2MG}{r}\right),
\]
with \(GM\) a time-independent integration constant, and \(f(t)\) an unknown
function of \(t\). The function \(f(t)\) can be made to equal unity by
defining a new time coordinate:
\[
  t'=\int^t f^{1/2}(t)\,\dd t.
\]
The metric is now entirely time-independent, and agrees with the Schwarzschild
solution \eqref{eq:8.2.12}. We have thus proved the \emph{Birkhoff
theorem},\hyperref[ch11-ref-30]{\textsuperscript{30}} that a spherically
symmetric gravitational field in empty space must be static, with a metric
given by the Schwarzschild solution.

The Birkhoff theorem is analogous to the result proved by Newton in his theory
of the lunar motion, that the gravitational field outside a spherically
symmetric body behaves as if the whole mass of the body were concentrated at
the center. It is a little surprising that this result should apply in general
relativity as well as in Newton's theory, for in general relativity a nonstatic
body will usually radiate gravitational waves. The Birkhoff theorem tells us
that, although a pulsating spherically symmetric body can of course produce
nonstatic gravitational fields within its mass, no gravitational radiation
can escape into empty space. In this sense, the Birkhoff theorem is analogous
to the well-known result of atomic theory, that a photon cannot be emitted in
a quantum transition between two states of zero spin.

The Birkhoff theorem may be applied, not only to the gravitational field
outside a body, but also to the field inside an empty spherical cavity at the
center of a spherically symmetric (but not necessarily static) body. In this
case the metric is again given by the Schwarzschild solution, but since the
point \(r=0\) is here in empty space, there can be no singularity, so the
integration constant \(MG\) must vanish. The Birkhoff theorem thus has the
corollary that \emph{the metric inside an empty spherical cavity at the center
of a spherically symmetric system must be equivalent to the flat-space
Minkowski metric \(\eta_{\mu\nu}\)}. This corollary is analogous to another
famous result of Newtonian theory, that the gravitational field of a spherical
shell vanishes inside the shell. Stars do not usually have holes at their
centers, so this corollary will not be of much use to us in this chapter. Its
importance arises from the fact that the Birkhoff theorem is a local theorem,
not depending on any conditions on the metric for \(r\to\infty\) (aside from
spherical symmetry), so that space must be flat in a spherical cavity at the
center of a spherically symmetric system, even if the system is
infinite---even, in fact, if the system is the whole universe. We shall see
in Section~15.1 that the corollary to Birkhoff's theorem can be used to justify
a limited use of Newtonian mechanics in cosmological problems.
